\documentclass{article}

% NeurIPS 2026 submission style. For a preprint, switch to
% \usepackage[preprint]{neurips_2026}; for camera-ready, add `final'.
\usepackage{neurips_2026}

\usepackage[utf8]{inputenc}
\usepackage[T1]{fontenc}
\usepackage{hyperref}
\usepackage{url}
\usepackage{booktabs}
\usepackage{amsmath}
\usepackage{amsfonts}
\usepackage{nicefrac}
\usepackage{microtype}
\usepackage[table]{xcolor}
\usepackage{multirow}
\usepackage{graphicx}
\usepackage{fvextra}          % extends fancyvrb with working breaklines
\usepackage{algorithm}
\usepackage{algorithmic}
\usepackage{xspace}
\usepackage{enumitem}

% Method-name macros. The current name is a placeholder; redefine here
% to rename throughout the paper.
\newcommand{\methodfull}{Agentic Conversation Context Curation\xspace}
\newcommand{\method}{AC3\xspace}

% Model logo + abbreviation macros for the multi-model megatable.
% Logos are normalized to 200x200 px upstream so heights match here.
\newcommand{\logogpt}{\raisebox{-0.20ex}{\includegraphics[height=1.7ex]{assets/logos/openai.pdf}}}
\newcommand{\logodsv}{\raisebox{-0.20ex}{\includegraphics[height=1.7ex]{assets/logos/deepseek.pdf}}}
\newcommand{\logokimi}{\raisebox{-0.20ex}{\includegraphics[height=1.7ex]{assets/logos/moonshot.png}}}
\newcommand{\mgpt}{\logogpt\,gpt-5.4}
\newcommand{\mdsv}{\logodsv\,dsv4f}
\newcommand{\mkimi}{\logokimi\,k2.6}

% A command to typeset a prompt from a file, with wrapping
\newcounter{promptctr}
\newcommand{\promptinput}[2][]{%
  \VerbatimInput[
    fontsize=\footnotesize,
    breaklines=true,
    breakanywhere=true,
    breaksymbolright=\small$\hookrightarrow$,
    obeytabs=true,
    tabsize=2,
    frame=leftline,
    framerule=0.6pt,
    rulecolor=\color{gray},
    #1
  ]{#2}%
}

\definecolor{darkblue}{rgb}{0, 0, 0.5}
\hypersetup{colorlinks=true, citecolor=darkblue, linkcolor=darkblue, urlcolor=darkblue}


\title{Agentic Context Management for Multi-Turn Human--AI Conversations}

% The NeurIPS submission style anonymizes authors automatically;
% the author block below is a placeholder for camera-ready.
\author{%
  Anonymous Authors
}

\newcommand{\fix}{\marginpar{FIX}}
\newcommand{\new}{\marginpar{NEW}}

% Custom 7-step gradient interpolated from base #56949f
% (generated by misc/colour_gradient_gen.py with lightness 0.92 -> 0.50;
%  softer dark end than the default 0.30).
\definecolor{tshade0}{HTML}{E0F1F4}
\definecolor{tshade1}{HTML}{C6E5EB}
\definecolor{tshade2}{HTML}{ABD9E1}
\definecolor{tshade3}{HTML}{91CDD8}
\definecolor{tshade4}{HTML}{77C2CF}
\definecolor{tshade5}{HTML}{5CB6C6}
\definecolor{tshade6}{HTML}{42AABC}

% Gradient-shaded score cell. Usage: \score{NN} -> shaded cell with NN.
% The trailing percent symbol is omitted: column headers should make units explicit.
% Bin boundaries chosen so 50 (chance on win-rate) sits mid-scale.
\newcommand{\score}[1]{%
\ifnum#1<20 \cellcolor{tshade0}\else%
\ifnum#1<35 \cellcolor{tshade1}\else%
\ifnum#1<50 \cellcolor{tshade2}\else%
\ifnum#1<65 \cellcolor{tshade3}\else%
\ifnum#1<78 \cellcolor{tshade4}\else%
\ifnum#1<88 \cellcolor{tshade5}\else%
\cellcolor{tshade6}\fi\fi\fi\fi\fi\fi%
#1}
% Variant: \scored{INT-FOR-COLOR}{DISPLAY} - integer for color bin, literal display string.
\newcommand{\scored}[2]{%
\ifnum#1<20 \cellcolor{tshade0}\else%
\ifnum#1<35 \cellcolor{tshade1}\else%
\ifnum#1<50 \cellcolor{tshade2}\else%
\ifnum#1<65 \cellcolor{tshade3}\else%
\ifnum#1<78 \cellcolor{tshade4}\else%
\ifnum#1<88 \cellcolor{tshade5}\else%
\cellcolor{tshade6}\fi\fi\fi\fi\fi\fi%
#2}

\begin{document}

\maketitle

\begin{abstract}
A growing body of work has identified a systematic failure mode in multi-turn LLM interaction: the model's own earlier outputs accumulate in context and bias later turns. Early incorrect inferences may anchor LLM reasoning even after the user supplies correcting information, a phenomenon called \emph{context pollution}. Previous proposed interventions discard assistant messages and rely on user turns containing sufficient information. However, more complex interactive applications are often \emph{referential}: partial work, prior decisions, and environmental state may only exist in assistant turns, but referenced in later turns by the user and assistant alike. We propose \textbf{\methodfull{} (\method)}, a training-free, inference-time method in which a separate \emph{analyzer} (an agentic subroutine) audits each turn in three steps: (1) consolidate a clean, unified specification from the user's messages alone, (2) compare the assistant's outputs against that specification, and (3) edit the context to keep verified work and remove invalidated reasoning. A non-obvious finding falls out of the design: pollution is \emph{contagious}. If the analyzer is shown the assistant's prior reasoning while consolidating the specification, downstream accuracy drops \emph{below} doing nothing, motivating context management for the context-management subagents themselves. Across four benchmarks of increasing referentiality (sharded single-turn questions in LiC, conversational tasks in CollabLLM, real human--AI dialogues in WildChat, and stateful agentic tool use in tau2-bench), \method is the only approach robust across the spectrum: it closes 55--80\% of the multi-turn gap on self-contained tasks (and exceeds the design-oracle on database), wins 84--86\% pairwise on real human--AI conversations, and remains viable in agentic tool use where blanket omission collapses from $\sim$60\% to 0\%. Gains generalize across four model families at different scales; selective curation, not blanket removal, is what scales to multi-turn human--AI interaction.
\end{abstract}


%% ============================================================
\section{Introduction}
\label{sec:intro}
%% ============================================================

\begin{figure}[t]
\centering
\includegraphics[width=\textwidth]{assets/ctxe_story.drawio.png}
\caption{\textbf{Selective context curation outperforms blanket removal as task referentiality grows.} Schematic of qualitative trends from Tables~\ref{tab:main}--\ref{tab:wildchat}; benchmarks are ordered by how much of the assistant's prior turns the next user turn references (LiC: self-contained sharded queries; CollabLLM: conversational back-and-forth; WildChat: real human--AI dialogues; tau2-bench: stateful tool-use, with results living only in assistant turns). \emph{Full context} (vanilla) plateaus throughout. \emph{Assistant Omission (AO)} is near-optimal at the self-contained end but degrades as discarded turns carry more task-relevant content, collapsing to 0\% on tau2-bench. \emph{\methodfull{} (\method, ours)} substantially outperforms vanilla on the self-contained and conversational ends, stays within trial noise of vanilla on tau2-bench where AO collapses entirely, and so remains viable across the spectrum: blanket removal is not a general fix for multi-turn degradation, but selective editing is.}
\label{fig:story}
\end{figure}

Multi-turn interaction is the dominant mode of use for large language models, from coding assistants that debug across dozens of exchanges to general-purpose computer-use agents like OpenClaw~\citep{openclaw2026}. Yet LLMs perform substantially worse in multi-turn settings than on equivalent single-turn tasks: \citet{laban2025lost} showed that accuracy drops 39\% on average across 15 models and 6 task domains, even in interactions spanning only $\sim$8 turns. They also identified the self-reinforcing dynamic that drives the gap: early errors about the user's intent get recorded in the assistant's own messages and bias later turns, so errors compound rather than self-correct as new user information arrives.
\citet{huang2026context} explores this \emph{context pollution} phenomenon in real human-AI conversations through a coarse but effective fix: \emph{assistant omission} (AO), which discards all assistant messages from context.

AO works precisely when user utterances by themselves contain sufficient information for task completion, in other words: \emph{self-contained}.
Many real interactions are \emph{referential} instead.
In tool-use settings where the agent manipulates environmental state, prior decisions and state tracking may live only in assistant turns.
Even in plain chat settings, users routinely refer to and build on prior assistant outputs (e.g.\ ``sharpen sentence~2''; ``add a column to that query'').
In these cases, removing the assistant messages leaves subsequent turns uninterpretable by deleting the content they refer to.

We propose \textbf{\methodfull{} (\method)}, a training-free, inference-time method that runs inside the agent harness before each assistant turn (Figure~\ref{fig:method}). Between turns, the harness invokes an \emph{agentic analyzer pipeline} that consolidates a clean, unified task specification from the user's messages, audits the assistant's prior outputs against that specification, and then edits the conversation context (by augmenting it with the analysis, resetting it from the clean specification, or rewriting it) to keep verified work and remove invalidated reasoning. The approach extends test-time scaling~\citep{wei2022chain,wang2023selfconsistency,yao2024tree,snell2024scalettc} to the multi-turn setting: rather than searching over outputs, we invest additional inference compute to curate the input context the model reasons over. This differs from prior context-management work~\citep[][\emph{inter alia}]{schroeder2025thread,zhang2025rlm,sun2025contextfolding} in two ways. (1) Prior methods address an agent's own thinking and execution trajectory following a single user request, while we address ongoing multi-turn user--agent dialogue. (2) they compact tokens for efficient long-horizon execution, while we remove pollution to improve correctness (Section~\ref{sec:related}).

This design is also shaped by an unexpected finding: pollution is \emph{contagious}. Even a separate model tasked with auditing the conversation anchors on the assistant's prior reasoning if exposed to it; pipelines that do not account for this perform \emph{worse} than no intervention at all. Context management must therefore extend to the analyzer subagents themselves: each is given only the conversation slices its role requires. Simply prompting the model to ignore content does not suffice (Section~\ref{sec:cognitive-hazard}).

Across four benchmarks of increasing referentiality, summarized in Figure~\ref{fig:story}, \method is the only method that holds up across the spectrum (Tables~\ref{tab:main},~\ref{tab:megatable},~\ref{tab:wildchat}): it closes 55--80\% of the multi-turn gap (and exceeds the oracle on database) on self-contained LiC tasks~\citep{laban2025lost} using Gated Reset as a single operator (e.g., 60.0\%$\to$80.0\% on math, 15.8\%$\to$64.4\% on code; mean over $N{=}3$ replay-mode reruns), wins 84--86\% pairwise on real human--AI conversations from WildChat~\citep{zhao2024wildchat}, and remains viable on stateful agentic tool-use tasks from tau2-bench~\citep{barres2025tau2}, where AO catastrophically fails (60.0\%$\to$0\%) by destroying tool-call results that live only in assistant turns. These gains generalize across four model families at different scales (closed and open-weight; +20--42pp average), and memory-based learning further improves without parameter updates.



%% ============================================================
\section{Problem Setting}
\label{sec:problem-setting}
%% ============================================================

\begin{figure}[t]
\centering
\includegraphics[width=\textwidth]{assets/ctxe_problem.png}
\caption{Context pollution in multi-turn conversations. The user reveals the task incrementally across turns; an early incorrect interpretation by the assistant persists in the context window and biases the model's generation even after the user provides correcting information. The deficit is behavioral, not capability-based: the same model solves the task when its own flawed prior reasoning does not dominate attention. Our method (Section~\ref{sec:methods}, Figure~\ref{fig:method}) supplies the missing selective attention at the harness level by rewriting the conversation to keep verified work and remove invalidated reasoning.}
\label{fig:problem}
\end{figure}

\subsection{Setup and notation}
Consider a multi-turn interaction. At turn $t$, the conversation context is the message sequence $C_t = (u_1, a_1, u_2, a_2, \ldots, u_t)$, where $u_i$ and $a_i$ denote user and assistant messages respectively, and the assistant samples $a_t \sim p_\theta(\cdot \mid C_t)$. Let $U_t = (u_1, \ldots, u_t)$ denote the user prefix and $A_{t-1} = (a_1, \ldots, a_{t-1})$ the assistant prefix.

\subsection{Context pollution as distribution shift}
Let $\mathcal{T}^*$ denote the true task specification, which the user reveals incrementally across $U_t$.  In principle, $U_t$ provides a progressively more complete view of $\mathcal{T}^*$.  However, $A_{t-1}$ encodes latent assumptions $\hat{\mathcal{T}}_i$ that may diverge from $\mathcal{T}^*$ as new user messages refine or contradict earlier interpretations.  This creates a \emph{distribution shift}: the model conditions on $C_t = (U_t, A_{t-1})$, but $A_{t-1}$ biases the posterior toward stale assumptions, so that
\begin{equation}
p_\theta(\text{correct} \mid U_t, A_{t-1}) < p_\theta(\text{correct} \mid U_t).
\end{equation}
The assistant's own prior outputs \emph{reduce} the probability of a correct response. Figure~\ref{fig:problem} illustrates the dynamic on a concrete trajectory: an early misinterpretation persists in the context window and re-anchors the model's reasoning even after the user has provided enough information to correct it. \citet{huang2026context} empirically confirm the inequality: simply omitting $A_{t-1}$ (assistant omission, AO) recovers much of the lost accuracy.

\subsection{Why omission is not enough: referential vs.\ self-contained turns}
\label{sec:referential}
AO depends on the assumption that $U_t$ alone fully determines a correct $a_t$. We call such turns \emph{self-contained}; their counterpart is \emph{referential}, where $a_t$ depends on content that exists only in $A_{t-1}$. Referential turns are pervasive. Assistant turns may contain previous decision making, partial artifacts (e.g., writing/programming drafts), or other forms of state information. For example, tool-call results are often recorded in $A_{t-1}$; omitting them causes the agent to repeat actions in the best case, and get totally lost in the environment in the worst case. Even in simpler chat applications without environmental state besides the conversation text, users build on assistant-generated content (editing a draft the assistant wrote, refining code from a previous turn, asking a follow-up that points at a specific earlier step~\citep{huang2026context}), so discarding those turns breaks the reference.
% Real conversations are typically referential; sharded LiC~\citep{laban2025lost} sits at the self-contained extreme and stateful tool-use~\citep{barres2025tau2} at the other.\footnote{We adopt the user-utterance-level terminology of \citet{choi2021decontextualization} (decontextualization). The closely-related ``stateless''/``stateful'' framing also fits, but conflates conversational reference (about prior \emph{utterances}) with environment state (about the world). We reserve ``stateful'' for the latter sense (e.g., tau2-bench tool calls), and use ``referential'' as the broader property that also covers plain chat.}

\subsection{Editing operator desiderata}
We seek an \emph{editing operator} $\mathcal{E}$ that produces a modified context $C'_t = \mathcal{E}(C_t)$ satisfying two properties:
\begin{enumerate}
    \item \textbf{Decontamination.} $C'_t$ suppresses the biasing effect of stale assumptions in $A_{t-1}$, so that $p_\theta(\text{correct} \mid C'_t) \geq p_\theta(\text{correct} \mid C_t)$.
    \item \textbf{Preservation.} $C'_t$ retains task-relevant information from both $U_t$ and $A_{t-1}$, so that $p_\theta(\text{correct} \mid C'_t) \geq p_\theta(\text{correct} \mid U_t)$ whenever $A_{t-1}$ contains useful referential state.
\end{enumerate}
AO satisfies the first property by construction (replacing $A_{t-1}$ with $\emptyset$) but trivially violates the second: it cannot retain what it discards. We need both.


%% ============================================================
\section{Methods}
\label{sec:methods}
%% ============================================================

\begin{figure}[t]
\centering
\includegraphics[width=\textwidth]{assets/ctxe_method(1).png}
\caption{\methodfull{} (\method). A separate analyzer audits the conversation in two subtasks before the assistant generates: Subtask~1 (spec extraction) sees only user turns, so it cannot be biased by the assistant's prior reasoning; Subtask~2 (approach evaluation) compares the clean spec against the assistant's full history. An intervention operator (Augment, Reset, or Rewrite) then edits the context. Structural exclusion in Subtask~1 is essential: relying on prompt-level instructions alone to ignore assistant turns collapses performance below baseline (Table~\ref{tab:cognitive-hazard}).}
\label{fig:method}
\end{figure}

We implement the editing operator $\mathcal{E}$ from Section~\ref{sec:problem-setting} as an agentic pipeline that analyzes the conversation history and edits the context before the next assistant turn. The core challenge is that an LLM cannot distinguish its own correct work from its own mistakes when both are encoded in the same assistant messages and influence future generation equally. Our method addresses this by externalizing the judgment: a separate analyzer consolidates what the user actually asked for (from user messages alone, so the consolidation is not biased by the assistant's framing), compares the assistant's work against that clean specification, and rewrites the conversation to keep what is correct and remove what is harmful. Algorithm~\ref{alg:context-edit} gives the full procedure.

\subsection{Conversation analysis via subtask decomposition}
\label{sec:decomposition}

Implementing $\mathcal{E}$ requires answering a prior question: \emph{how do you identify context pollution?} Pollution is content that biases the model toward an incorrect approach, but ``incorrect'' is only meaningful against ground truth. In multi-turn conversation, the user's messages are the primary source of truth: an approach is invalidated when it rests on assumptions contradicted by later user messages. Identifying pollution is hard because new user information can look compatible with a polluted approach when it actually requires a different one. Users may also revise earlier requirements across turns, so the spec must be \emph{consolidated} from the union of user messages to resolve this.

These observations yield two design pressures. First, identifying pollution requires a task specification \emph{grounded} in user requirements, not derived from the assistant's interpretation; this aligns with prior findings that LLMs cannot reliably self-correct without external feedback~\citep{kamoi2024criticalselfcorrect} and cannot locate their own reasoning errors but can correct them given the error location~\citep{tyen2024errors}. Our user-grounded spec provides that external anchor. Second, extraction itself must be \emph{shielded} from the assistant's reasoning, lest it fall victim to pollution. We therefore decompose the analysis into two subtasks (Algorithm~\ref{alg:context-edit}, lines 2--5; tested in Section~\ref{sec:cognitive-hazard}):
\begin{itemize}[leftmargin=*, topsep=2pt, itemsep=1pt, parsep=0pt]
    \item \textbf{Subtask 1: spec extraction (selective attention).} The model sees only $U_t$ (user + system messages), with $A_{t-1}$ structurally excluded, and produces a consolidated task spec $\hat{\mathcal{T}}$. The system message grounds the output format (e.g., SQL schema, function signatures). Physically excluding assistant messages implements attentional control the model cannot perform internally.
    \item \textbf{Subtask 2: approach evaluation (error monitoring).} The model receives $\hat{\mathcal{T}}$, the full context $C_t$ (including assistant messages), and optionally a memory cheatsheet $M$. Grounded in $\hat{\mathcal{T}}$, it identifies what the assistant got right (\texttt{aligned}) and where it diverged (\texttt{issues}).
\end{itemize}

\subsection{Intervention strategies}
\label{sec:strategies}

The analysis yields $(\hat{\mathcal{T}}, \texttt{aligned}, \texttt{issues})$. Three operators implement $\mathcal{E}$ with increasing context modification; gating is orthogonal. The analytical core (Section~\ref{sec:decomposition}) is the contribution and operator choice is a design knob (defaults per task in Table~\ref{tab:main}).
\begin{itemize}[leftmargin=*, topsep=2pt, itemsep=1pt, parsep=0pt]
    \item \textbf{Augment.} Append the artifact: $C'_t = C_t \oplus (\hat{\mathcal{T}}, \texttt{aligned}, \texttt{issues})$. Adds corrective signal but leaves polluting content in $A_{t-1}$ visible, where it may still anchor generation.
    \item \textbf{Reset.} Populate a template that \emph{replaces} the conversation: $C'_t = \textsc{Template}(\hat{\mathcal{T}}, \texttt{aligned})$. Because the analysis maps deterministically (no extra LLM call), this is the cheapest intervention with explicit removal and avoids the summarization losses Rewrite can introduce (Section~\ref{sec:lic-results}).
    \item \textbf{Rewrite.} An additional LLM call generates a restructured context: $C'_t \sim p_\theta(\cdot \mid \hat{\mathcal{T}}, \texttt{aligned}, \texttt{issues}, C_t)$. Most flexible; trades occasional summarization loss for arbitrary restructurings, with a ceiling we expect to surpass deterministic mapping over time.
    \item \textbf{Gating} (orthogonal). Any operator can be applied unconditionally or \emph{gated}, triggered only on flagged issues (Algorithm~\ref{alg:context-edit}, lines 6--8). Gating reduces cost and avoids disrupting on-track conversations --- critical in agentic settings where resets discard expensive-to-reobtain tool results.
\end{itemize}
% Full pseudocode is in Algorithm~\ref{alg:context-edit} (Appendix~\ref{app:algorithm}).

\subsection{Memory-based learning}
\label{sec:memory}

We explore whether context curation can be learned from experience via lightweight memory. 
To isolate the impact of memory-based learning from any advanced implementation details, we adapt the Dynamic Cheatsheet~\citep{suzgun2025cheatsheet}, a minimal design that represents memory as a single text artifact updated globally without a retrieval mechanism (removing a potential confounding factor).

After each problem instance, a reflection query examines the analysis and outcome, abstracting transferable takeaways (e.g., common pitfalls, effective consolidation strategies). These update a cheatsheet $M$, a plain-text memory of strategies and pitfalls accumulated across instances. The cheatsheet is injected into the analyzer's approach evaluation (Subtask~2) on subsequent problems, providing learned priors that improve the quality of $\mathcal{E}$. We cap cheatsheet length at 1500 words to prevent unbounded growth from degrading quality. Implementation details on update strategies and evaluation protocols are in Appendix~\ref{app:memory-details}.


%% ============================================================
\section{Experiments}
\label{sec:experiments}
%% ============================================================

Unless otherwise noted, GPT-5-mini serves as both assistant and analyzer throughout. The four benchmarks below differ in how much of the assistant's prior work is essential to task completion, which lets us read how each intervention responds to increasing referentiality.

\noindent\textbf{Lost in Conversation (LiC).}
LiC~\citep{laban2025lost} shards a fully-specified single-turn question into fragments revealed incrementally by a simulated user. We test four domains: math, code, database (SQL), and function calling (``actions''). We compare against Assistant Messages Omitted (AO;~\citealt{huang2026context}), the Concatenate User baseline~\citep{laban2025lost} (joins the shards back into a single message), and ERGO~\citep{khalid2025ergo} (resets via LLM rewrite of prior user messages). Because each LiC instance is a single-turn question by construction, AO and Concat User are \emph{design-oracle} baselines. Replay-mode protocol and per-task adjustments are in Appendix~\ref{app:user-sim-qc}.

\noindent\textbf{CollabLLM. \citep{wu2025collabllm}}
We evaluate on MATH-Hard and BigCodeBench using the CollabLLM multi-turn framework. Unlike LiC, the conversation beats are not pre-planned in a sharded format. Instead the simulated user lazily issues a partial request (still corresponding to an originally single-turn benchmark) and both sides exchange natural follow-ups, so the assistant's prior turns contain partial work that may matter downstream. Evaluation details are in Appendix~\ref{app:collabllm}.

\noindent\textbf{WildChat.}
We evaluate on WildChat-1M~\citep{zhao2024wildchat} following \citet{huang2026context}: real human--assistant dialogue where users reference and build on earlier assistant outputs, so AO can sever the conversational thread and is no longer an oracle. Filtering, two-phase setup, and pairwise judge are in Appendix~\ref{app:wildchat-turns}; Appendix~\ref{app:trajectories} walks through a Maven/Java debugging conversation.

\noindent\textbf{Tau2-bench.}
As a preliminary exploration of the agentic setting, we evaluate on the \texttt{telecom\_small} subset of tau2-bench~\citep{yao2024tau}, simulating a customer-service agent that resolves a telecom troubleshooting call via diagnostic tool calls. This is a setting with explicit environmental state. Unlike other agentic benchmarks, tau2-bench makes conversation especially meaningful with a novel dual-control setup where user and agent are both capable of manipulating state with unique capabilities. This forms a setting where state information organically lives across both parties' messages and actions, and thereby serves as a stress test for context management in the face of complex state tracking. Agentic adaptation details are in Appendix~\ref{app:agentic-methods}.

%% ============================================================
\section{Results}
\label{sec:results}
%% ============================================================

\subsection{Lost in Conversation}
\label{sec:lic-results}

\begin{table}[!htbp]
\centering
\caption{\textbf{Per-task LiC results (GPT-5-mini, accuracy \%).} Cells shaded by score (darker $=$ higher). Indented \texttt{+ Memory} rows add the memory-based learner. $^{\diamond}$ marks \emph{design-oracle} baselines (collapsing multi-turn to single-turn by construction); \method \emph{exceeds} them in two cells marked $^{\dagger}$. $^{\ddagger}$ Gated-Reset row reports the mean over $N{=}3$ replay-mode reruns per cell (last-turn regen with fixed user-sim trajectory); other ours-rows are single-trial point estimates --- see Appendix~\ref{app:variance}. Multi-respondent results for LiC, CollabLLM, and tau2-bench are in Table~\ref{tab:megatable}; multi-respondent WildChat results are in Table~\ref{tab:wildchat}.}
\label{tab:main}

\begin{tabular}{lcccc}
\toprule
\textbf{Strategy} & \textbf{Math}~$\uparrow$ & \textbf{Code}~$\uparrow$ & \textbf{Database}~$\uparrow$ & \textbf{Actions}~$\uparrow$ \\
\midrule
\multicolumn{5}{l}{\emph{Prior work \& baselines}} \\
Baseline (full context) & \scored{60}{60.0} & \scored{15}{15.8} & \scored{4}{4.0} & \scored{34}{34.8} \\
\quad + Memory & \scored{55}{55.0} & \scored{21}{21.1} & \scored{4}{4.0} & \scored{34}{34.8} \\
AO~\citep{huang2026context}$^{\diamond}$ & \scored{85}{85.0} & \scored{77}{77.8} & \scored{32}{32.0} & \scored{82}{82.6} \\
Concat User~\citep{laban2025lost}$^{\diamond}$ & \scored{84}{84.2} & \scored{68}{68.4} & \scored{32}{32.0} & \scored{87}{87.0} \\
ERGO~\citep{khalid2025ergo} & \scored{69}{69.6} & \scored{44}{44.0} & \scored{12}{12.0} & \scored{48}{48.0} \\
\midrule
\multicolumn{5}{l}{\emph{\methodfull{} (ours)}} \\
\method-Augment & \scored{80}{80.0} & \scored{55}{55.6} & \scored{32}{32.0} & \scored{47}{47.8} \\
\quad + Memory & \textbf{\scored{90}{90.0}}$^{\dagger}$ & \scored{68}{68.4} & \scored{44}{44.0} & \scored{47}{47.8} \\
\method-Reset & \scored{75}{75.0} & \scored{61}{61.1} & \textbf{\scored{48}{48.0}}$^{\dagger}$ & \scored{52}{52.2} \\
\quad + Memory & \scored{85}{85.0} & \scored{68}{68.4} & \scored{44}{44.0} & \scored{52}{52.2} \\
\method-Gated-Reset$^{\ddagger}$ & \scored{80}{80.0} & \scored{64}{64.4} & \scored{39}{38.7} & \scored{61}{61.3} \\
\bottomrule
\end{tabular}
\end{table}

\begin{table}[t]
\centering
\caption{\textbf{Multi-respondent robustness on LiC, CollabLLM, and tau2-bench.} Each cell is benchmark-averaged accuracy (LiC, CollabLLM) or task reward (tau2-bench, \texttt{telecom\_small}, $n{=}19$ per cell) for one (strategy, respondent) pair. Respondents: \mgpt~(Azure OpenAI), \mdsv~(DeepSeek-V4-Flash, Azure Foundry), \mkimi~(Kimi-K2.6, Azure Foundry). The \method-Reset row reports always-on Reset where we ran it; cells marked $^{*}$ are Gated-Reset (tau2-bench, where always-on Reset was not run). \emph{LiC} averages math/code/database/actions accuracy (per-task GPT-5-mini view in Table~\ref{tab:main}). \emph{CollabLLM} averages MATH-Hard and BigCodeBench; cells marked $^{\dagger}$ report MATH-Hard only because the BigCodeBench Reset cell hit Azure content-filter rejections (gpt-5.4) or Foundry endpoint instability (Kimi). \mkimi~tau2-bench Baseline/AO cells are rate-limit-clipped floors at Foundry workers=4 (true Baseline likely 40--50\%; AC3 cells ran clean). WildChat is reported separately in Table~\ref{tab:wildchat}.}
\label{tab:megatable}
\footnotesize
\setlength{\tabcolsep}{4pt}
\begin{tabular}{l ccc | ccc | ccc}
\toprule
& \multicolumn{3}{c|}{\textbf{LiC}} & \multicolumn{3}{c|}{\textbf{CollabLLM}} & \multicolumn{3}{c}{\textbf{tau2-bench}} \\
\textbf{Strategy} & \mgpt & \mdsv & \mkimi & \mgpt & \mdsv & \mkimi & \mgpt & \mdsv & \mkimi \\
\midrule
Baseline                 & \scored{60}{60.4} & \scored{51}{51.3} & \scored{60}{60.4} & \scored{55}{55.0} & \scored{50}{50.0} & \scored{57}{57.5} & \scored{68}{68.4} & \scored{31}{31.6} & \scored{26}{26.3} \\
AO                       & \scored{74}{74.5} & \scored{69}{69.5} & \scored{68}{68.2} & \scored{57}{57.5} & \scored{52}{52.5} & \scored{63}{63.2} & \scored{0}{0.0}   & \scored{0}{0.0}   & \scored{0}{0.0}   \\
\method-Augment          & \scored{77}{77.0} & \scored{67}{67.0} & \scored{73}{73.8} & \scored{56}{56.3} & \scored{57}{57.5} & \scored{55}{55.4} & \scored{84}{84.2} & \scored{57}{57.9} & \scored{57}{57.9} \\
\method-Reset            & \scored{77}{77.2} & \scored{68}{68.4} & \scored{74}{74.3} & \scored{47}{47.0$^{\dagger}$} & \scored{52}{52.5} & \scored{47}{47.0$^{\dagger}$} & \scored{52}{52.6$^{*}$} & \scored{47}{47.4$^{*}$} & \scored{68}{68.4$^{*}$} \\
\method-Rewrite          & \scored{76}{76.6} & \scored{64}{64.7} & \scored{74}{74.8} & \scored{53}{53.8} & \scored{44}{44.9} & \scored{58}{58.3} & \scored{73}{73.7} & \scored{57}{57.9} & \scored{73}{73.7} \\
\bottomrule
\end{tabular}
\end{table}

\begin{table}[t]
\centering
\caption{\textbf{WildChat: pairwise win rates per respondent (Phase-2 turns).} Each cell is a GPT-5-mini-judged head-to-head quality win-rate against the baseline indicated (AO = assistant omission; FC = full context). Higher is better; 50\% is parity. The GPT-5-mini column shows the original Phase-2 numbers (Table~\ref{tab:wildchat-memory} provides the corresponding DeepSeek memory ablation); the three Foundry columns use the Phase-1-prefix evaluation set ($n{\approx}48{-}74$ per cell, gpt-5-mini analyzer locked across respondents to isolate the rewriter-prompt variable). The \mgpt~Gated-Reset cell ($-$14.5pp vs.\ always-on Reset on the same prefixes) is discussed in Section~\ref{sec:wildchat-results}; corresponding cells on the Foundry respondents were not run.}
\label{tab:wildchat}
\footnotesize
\setlength{\tabcolsep}{4pt}
\begin{tabular}{l cc | cc | cc | cc}
\toprule
& \multicolumn{2}{c|}{\textbf{GPT-5-mini}} & \multicolumn{2}{c|}{\mgpt} & \multicolumn{2}{c|}{\mdsv} & \multicolumn{2}{c}{\mkimi} \\
\textbf{Strategy} & vs AO & vs FC & vs AO & vs FC & vs AO & vs FC & vs AO & vs FC \\
\midrule
\method-Augment      & ---              & ---              & \scored{84}{84.2} & \scored{77}{77.2} & \scored{84}{84.2} & ---              & \scored{85}{85.7} & ---              \\
\method-Reset        & \scored{83}{83.0} & \scored{83}{83.5} & \scored{88}{88.6} & \scored{77}{77.3} & \scored{75}{75.0} & ---              & \scored{71}{71.6} & ---              \\
\method-Gated-Reset  & \scored{86}{86.1} & \scored{83}{83.8} & \scored{74}{74.1} & \scored{72}{72.4} & ---              & ---              & ---              & ---              \\
\method-Rewrite      & \scored{82}{82.6} & \scored{80}{80.3} & \scored{83}{83.3} & \scored{72}{72.9} & \scored{79}{79.2} & \scored{73}{73.6} & \scored{91}{91.5} & \scored{76}{76.3} \\
\bottomrule
\end{tabular}
\end{table}

\noindent\textbf{Context editing closes most of the gap to the design oracle.}
Reading the Gated-Reset row (mean over $N{=}3$ replay-mode reruns; Table~\ref{tab:main}): on math, Gated Reset reaches 80.0\%, closing 80\% of the 25.0pp gap between baseline and AO; on code, Gated Reset reaches 64.4\%, closing 78\% of the 62.0pp gap. On database, Reset (48.0\%) \emph{exceeds} both oracle baselines (32.0\%) because consolidation outperforms na\"ive removal when shards need schema-grounded assembly. Per-run variance on the Gated-Reset row is in the 5--7pp range and the means cleanly clear both the Full-Context baseline and the no-memory Augment operator on every LiC task, so the headline LiC ordering is not a lucky-run artifact (Appendix~\ref{app:variance}).

\noindent\textbf{Context removal provides additional benefit.}
Reset and Gated Reset outperform Augment on code (+5.5--8.8pp) and database (+6.7--16.0pp); the difference (same analysis) is whether the assistant sees a polluted or clean context, confirming anchoring is separable from intent fragmentation.

\noindent\textbf{Memory amplifies analysis.}
Augment~+~Memory reaches 90.0\% on math, exceeding the AO oracle (85.0\%), and improves code by +12.8pp and database by +12.0pp; on Actions it is neutral, matching the no-memory Augment baseline at 47.8\%. Memory has no effect on Baseline, confirming memory improves analysis and edit quality, not raw task performance.

\noindent\textbf{Generalization across model families.}
\label{sec:multi-model}
We evaluate three additional models (GPT-5, DeepSeek V3.2 (37B active, open-weight), and Qwen 3.5 35B-A3B (3B active)) on a harder problem subset (Appendix~\ref{app:multi-model-protocol}). Reset improves accuracy in 15 of 16 model--task pairs (Table~\ref{tab:multi-model}), with gains of +20--42pp on average across closed and open-weight architectures at different scales. Tables~\ref{tab:megatable} and~\ref{tab:wildchat} extend this picture beyond GPT-5-mini: across three additional respondent models (GPT-5.4, DeepSeek-V4-Flash, Kimi-K2.6) and all four benchmarks, every \method~operator beats Baseline (full context) on every (model, benchmark) cell where we evaluated both, while AO collapses entirely on tau2-bench across all three. The qualitative ordering is benchmark-internal: on LiC, all three operators land within $\sim$1--5pp of each other; on tau2-bench the spread widens to 16--47pp and the winning operator changes by respondent (see Section~\ref{sec:tau2-results}).


\subsection{CollabLLM}
\label{sec:collab-results}

We use CollabLLM primarily to probe the structural-exclusion design choice in a setting with potential cross-turn referentiality. Because intent may weave across user and assistant turns, we run Rewrite (single-pass over the full transcript, no structural exclusion of assistant content) rather than Reset/Gated Reset, which assume a clean user-only specification. Rewrite improves over the full-context baseline on both tasks (Table~\ref{tab:megatable}, CollabLLM block; per-task GPT-5-mini numbers in Appendix~\ref{app:collabllm}), with a +20pp gain on BigCodeBench. AO is also competitive (50.0 / 82.5 vs.\ Rewrite 45.0 / 82.5 on GPT-5-mini): CollabLLM sits midway along the referentiality axis, where blanket omission still helps but is no longer dominant.
% The decisive separation between blanket removal and selective editing emerges further along the spectrum (\S\ref{sec:wildchat-results}, \S\ref{sec:tau2-results}), where AO actively hurts.


\subsection{WildChat}
\label{sec:wildchat-results}

Unlike LiC, WildChat conversations are referential, so AO is no longer an oracle and can actively degrade quality. Following \citet{huang2026context}, we classify each turn as \texttt{new\_ask} (the user introduces an unrelated topic), \texttt{feedback} (the user corrects or refines the assistant's prior response), or \texttt{no\_feedback} (the user continues without explicitly referencing the assistant's response, e.g., adding information or following up on their own prior message; see Appendix~\ref{app:wildchat-turns}). Phase~1 reproduces Huang's finding: AO helps on \texttt{new\_ask} and \texttt{feedback} turns but \emph{hurts} on referential \texttt{no\_feedback} turns, where FC wins the pairwise comparison.

\noindent\textbf{Context editing outperforms both AO and full context.}
Gated Reset achieves 86.1\% win rate vs.\ AO and 83.8\% vs.\ FC on GPT-5-mini while intervening on only 72.3\% of turns (Table~\ref{tab:wildchat}, GPT-5-mini column), with selective gating that intervenes less on referential \texttt{no\_feedback} turns than on \texttt{new\_ask} or \texttt{feedback} turns (Appendix~\ref{app:wildchat-turns}). DeepSeek V3.2 shows consistent gains (Reset wins 72.5\% vs.\ AO and 72.5\% vs.\ FC; Table~\ref{tab:wildchat-memory}). Appendix~\ref{app:trajectories} shows a Maven-debugging conversation where Gated Reset wins all 14 comparisons across 7 turns: AO re-suggests already-tried commands, FC anchors on the flawed diagnostic, and Gated Reset keeps the correct diagnosis while discarding broken reasoning.

\noindent\textbf{Gated vs.\ always-on Reset on a strong respondent.}
On most cells, Gated Reset and always-on Reset are interchangeable: the analyzer flags issues on $\ge$97\% of LiC and CollabLLM turns, so the gate almost never closes and the two strategies produce identical outputs. They diverge on WildChat with the strongest respondent (gpt-5.4): always-on Reset wins 88.6\% vs.\ AO (Table~\ref{tab:wildchat}), while Gated Reset on the same prefix set scores 74.1\%, a $-$14.5pp gap driven by false-negative \texttt{needs\_edit=False} decisions when the analyzer judges a turn ``on track'' but a clean reset would still have helped. The gating heuristic's failure mode is asymmetric: it is conservative about closing the gate when issues are visible, but its decision boundary is calibrated to weaker respondents and leaves quality on the table on stronger ones. We recommend always-on Reset when intervention cost is not a constraint and the respondent is near-ceiling; gated when interventions carry state-disruption cost (the tau2-bench regime, Section~\ref{sec:tau2-results}).

\noindent\textbf{Rewrite wins on referential generation.}
On WildChat with Kimi-K2.6, Rewrite wins 91.5\% vs.\ AO and substantially beats both Augment (85.7\%) and Reset (71.6\%) (Table~\ref{tab:wildchat}). Reset's template strips conversational structure that the next turn refers to; an open-ended Rewrite reorganizes the same analyzer output into a form that preserves what the user is pointing at. This is the cleanest evidence in the paper that flexibility (Rewrite) and rigidity (Reset) trade against each other along the referentiality axis.


\subsection{Tau2-bench}
\label{sec:tau2-results}

AO collapses to \textbf{0\%} on tau2-bench across every respondent we evaluated (Table~\ref{tab:megatable}, tau2 block): stripping assistant messages erases every diagnostic result the agent has fetched, inducing infinite re-lookup loops, regardless of the underlying respondent's capability; blanket omission is fundamentally incompatible with stateful tool use.

\noindent\textbf{The winning AC3 operator changes with respondent strength.} Across the three Foundry-served respondents (Table~\ref{tab:megatable}, tau2 panel), every AC3 operator beats Baseline, but the winner is model-dependent: gpt-5.4 (strongest) is best served by AC3-Augment (84.2\% vs.\ Baseline 68.4\%, $+$15.8pp) --- the analyzer's hints help without disrupting state; DeepSeek-V4-Flash is jointly best on Augment and Rewrite (both 57.9\% vs.\ Baseline 31.6\%, $+$26.3pp); Kimi-K2.6 (weakest base) is best served by the heaviest operator, AC3-Rewrite (73.7\% vs.\ Baseline 26.3\%, $+$47.4pp). The pattern matches the ``appropriate intensity'' framing: stronger base agents benefit from analyzer hints without destruction, weaker agents benefit from heavier interventions that aggressively prune polluted reasoning, and the same analyzer pipeline supports both extremes by composing different downstream operators.

\noindent\textbf{The gpt-5-mini cell sits at the edge of the curve.} On gpt-5-mini (Appendix~\ref{app:tau2-diagnostic}), Gated Reset stays within trial noise of the full-context Baseline across 3 seeds; only 1 of 11 baseline failures is attributable to context pollution (Table~\ref{tab:tau2-failure-modes}), and the customer-service policy's breadth-first diagnostic list partially substitutes for an external editor. This is consistent with the model-dependent pattern above: on respondents whose baseline failure mode is not pollution-driven, editing helps less; on respondents whose baseline failures \emph{are} pollution-driven (the Foundry sweep above), editing helps a great deal. Across all four respondents the load-bearing finding is the same: selective editing is the only context-management approach that remains viable when tool-call results live only in assistant turns.


%% ============================================================
\section{Discussion}
\label{sec:discussion}
%% ============================================================

\paragraph{Editing vs.\ removal across settings.}
The value of context curation over blanket removal scales with referentiality: on self-contained LiC, omission is a design oracle and we close most of the gap to it; on WildChat, editing outperforms both AO and full context (80--86\% win rates); on tau2-bench, AO collapses to 0\% on every respondent we tried while editing remains viable, with the Foundry sweep (Table~\ref{tab:megatable}) yielding +15.8 to +47.4pp over Baseline depending on the respondent. Gains generalize across model families (Tables~\ref{tab:multi-model} and~\ref{tab:megatable}). \method incurs heavier per-turn cost than detection-only alternatives such as \citet{huang2026context} (embedding classifier) and \citet{khalid2025ergo} (output-entropy spikes), but the analyzer's outputs are plain-text reasoning of the kind SOTA models already emit in their internal \texttt{<reasoning>} traces, so recent self-distillation methods~\citep{zhao2026selfdistillreason,hubotter2026rlselfdistill,shenfeld2026selfdistill} can absorb the analysis into the base model with no inference-time call overhead; embedding outputs and entropy spikes do not absorb as cleanly.

\paragraph{Appropriate intensity: one analyzer, three operators.}
A consistent pattern emerges from the multi-model results in Table~\ref{tab:megatable}: the right \emph{operator} depends on the respondent and the benchmark, but the analyzer pipeline that feeds them is shared. On tau2-bench, the strongest respondent (gpt-5.4) benefits most from the lightest operator (Augment, $+$15.8pp), while the weakest (Kimi-K2.6) benefits most from the heaviest (Rewrite, $+$47.4pp). On WildChat, where users explicitly refer back to assistant content, Rewrite's open-ended restructuring wins on Kimi-K2.6 ($+$5.8pp over Augment) while Reset's template wins on the strongest respondent (gpt-5.4) by leaving the analyzer-curated specification unedited by a second LLM call. This suggests presenting \method~as a pipeline that decouples analysis (one shared subroutine) from intervention intensity (a per-deployment knob), rather than as a single recommended operator.


\paragraph{Context pollution is contagious.}
\label{sec:analysis-ablations}
\label{sec:cognitive-hazard}
Letting the spec extractor see assistant messages drops accuracy \emph{below baseline}; even a model explicitly tasked as an independent reviewer anchors on the assistant's flawed reasoning. Counterintuitively, the contaminated single-pass variant then \emph{beats} the contaminated two-step version (55\% vs.\ 40\% on math) because the latter's downstream stage trusts a biased upstream spec as ground truth (Table~\ref{tab:cognitive-hazard}, Appendix~\ref{app:hard-attention}). Decomposition helps only when earlier stages produce clean outputs: a structural data-flow constraint that prompting alone cannot enforce.


\paragraph{Memory-based learning.}
\label{sec:memory-discussion}
Memory amplifies analysis when headroom exists: Augment~+~Memory reaches 90\% on LiC math (exceeding AO, Table~\ref{tab:main}) and improves DeepSeek's Reset on WildChat by +13.8pp (Table~\ref{tab:wildchat-memory}), while being neutral on near-ceiling GPT-5-mini. In a more realistic \emph{soft-attention} setting (Appendix~\ref{app:soft-attention}), memory partially closes the hard-attention gap but substantial headroom remains.




%% ============================================================
\section{Related Work}
\label{sec:related}
%% ============================================================

\noindent\textbf{Multi-turn degradation and context pollution.}
\citet{laban2025lost} introduced LiC; \citet{huang2026context} termed the failure mode \emph{context pollution} and showed AO recovers much of the lost performance; \citet{khalid2025ergo} similarly resets via entropy heuristics. All treat context as \emph{uniformly harmful}, an assumption that holds only for self-contained tasks (the LiC setting, where AO is a design-oracle); on referential and stateful settings, blanket omission collapses (Section~\ref{sec:tau2-results}). We instead treat context as \emph{heterogeneous}, distinguishing correct assistant work from polluting content.

\noindent\textbf{Agentic context management.}
Concurrent work on long-horizon tasks (THREAD~\citep{schroeder2025thread}, Recursive Language Models~\citep{zhang2025rlm}, Context Folding~\citep{sun2025contextfolding}, U-Fold~\citep{su2026ufold}, MemoBrain~\citep{qian2026memobrain}) targets \emph{compaction} (fitting more reasoning into bounded context); we address \emph{context pollution}, where prior outputs mislead generation within ample context. Detailed comparison and additional related work are in Appendices~\ref{app:additional-related}--\ref{app:agentic-context}.


%% ============================================================
\section{Conclusion}
\label{sec:conclusion}
%% ============================================================

We propose \methodfull{} (\method), a training-free, harness-level method that mitigates multi-turn degradation by selectively curating rather than discarding context. Across self-contained, referential, and stateful agentic regimes, \method outperforms both full context and blanket omission (Gated Reset closes 55--80\% of the LiC gap and exceeds the design oracle on database; 84--86\% pairwise win on WildChat; on tau2-bench, where AO collapses to 0\%, Gated Reset stays within trial noise of the full-context baseline rather than catastrophically failing). Gains generalize across four model families spanning closed and open-weight architectures (+20--42pp average), and a memory-based learner further improves analysis quality without parameter updates. A non-obvious finding falls out of the design: pollution is contagious, since even an independent reviewer model anchors on the assistant's reasoning if exposed to it, so the analysis pipeline must structurally exclude assistant turns rather than rely on prompting alone. As LLM interaction grows more referential and agentic, blanket removal cannot scale where selective curation can, and the gap between them will only widen.


%% ============================================================
\begin{ack}
%% ============================================================
% The \texttt{ack} environment is automatically hidden in the anonymous
% submission. Populate acknowledgments and funding disclosures for the
% camera-ready version.
\end{ack}


%% ============================================================
\bibliographystyle{plainnat}
\bibliography{neurips_2026_conference}


%% ============================================================
\appendix

\section{Algorithm}
\label{app:algorithm}

\begin{algorithm}[h]
\caption{\methodfull{} (\method)}
\label{alg:context-edit}
\begin{algorithmic}[1]
\REQUIRE Conversation context $C_t = (u_1, a_1, \ldots, u_t)$, gating flag $g$
\STATE \textbf{// Subtask 1: Selective attention (spec extraction)}
\STATE $U_t \leftarrow (u_1, \ldots, u_t)$ \hfill $\triangleright$ Structurally exclude $A_{t-1}$
\STATE $\hat{\mathcal{T}} \leftarrow \textsc{Analyze}(U_t)$ \hfill $\triangleright$ Extract clean task specification
\STATE \textbf{// Subtask 2: Error monitoring (approach evaluation)}
\STATE $(\texttt{aligned}, \texttt{issues}) \leftarrow \textsc{Evaluate}(\hat{\mathcal{T}}, C_t)$
\IF{$g$ \AND $\texttt{issues} = \emptyset$}
    \RETURN $C_t$ \hfill $\triangleright$ No edit needed (gated pass-through)
\ENDIF
\STATE \textbf{// Apply editing operator $\mathcal{E}$}
\STATE $C'_t \leftarrow \textsc{Intervene}(\hat{\mathcal{T}}, \texttt{aligned}, \texttt{issues}, C_t)$
\RETURN $C'_t$
\end{algorithmic}
\end{algorithm}


\section{Evaluation details}
\label{app:eval-details}

\subsection{LiC evaluation adjustments}
\label{app:user-sim-qc}

We make task-specific adjustments to the LiC evaluation protocol to reduce false negatives caused by overly strict extraction and ambiguous formatting, and apply a post-hoc false negative identification procedure across all tasks.

\paragraph{Replay-mode protocol.} To isolate intervention effects, we run all LiC strategies in \emph{replay mode}: every strategy shares the same baseline conversation prefix and only the final turns are regenerated under each strategy. Cell-to-cell comparisons therefore measure what an editor recovers from a fixed history rather than confounding strategy effects with divergent rollouts.

\paragraph{Sampling parameters.} All LLM calls (assistant, user-simulator, system-agent, and analyzer) use \texttt{temperature}~$=$~$1.0$ with default top-$p$. This is the reasoning-model default: GPT-5-family models silently override any \texttt{temperature=0.0} request to \texttt{1.0}, so we use \texttt{1.0} throughout for consistency. We do not pass a \texttt{reasoning\_effort} parameter for analyzer calls (constraining reasoning effort severely degrades editor quality on GPT-5-mini in our pilots). The same sampling configuration is used both when generating the original baseline trajectory and when each strategy regenerates the final turn(s) under replay mode.

\paragraph{API cost.} All experiments are inference-only and were run via commercial APIs (OpenAI, Anthropic) and open-weight inference endpoints (DeepSeek, Qwen). We report per-run costs rather than totals so the figures translate directly to reproduction budget. With \texttt{gpt-5-mini} as both assistant and analyzer, a single Gated-Reset replay-mode rerun of an LiC task ($n{=}18$--$25$ samples) costs approximately \$0.04--\$0.06; tau2-bench is more expensive per task because conversations are longer and tool-call payloads inflate token usage, at roughly \$0.70 per trial of the 20-task \texttt{telecom\_small} subset, with the unconditional Reset variant adding $\sim$+20\% over the Baseline due to the analyzer's 2--3 extra LLM calls per turn. WildChat and CollabLLM are scored by GPT-5 judge calls over pre-generated assistant responses, so per-experiment cost is dominated by judge tokens rather than full multi-turn simulation; a single judged WildChat condition over $n{\approx}175$ turns costs a few dollars. Despite our scaled intervention adding 2--3 LLM calls per turn, our project-level total is well under the $\sim$\$5{,}000 envelope reported by~\citet{laban2025lost}: \citeauthor{laban2025lost} ran six task domains across fifteen models with full multi-turn simulation, while our headline tables span four LiC tasks (with replay mode amortizing the baseline trajectories), three CollabLLM/WildChat conditions, one tau2-bench subset, and four model families.

\subsubsection{Task-specific prompt and extraction changes}
\label{app:eval-adjustments}

\paragraph{Math.} The original LiC math evaluator uses symbolic answer extraction that is overly strict, producing frequent false negatives when the assistant's final answer is correct but formatted unexpectedly. Since GSM8K problems have integer-only solutions, we simplify extraction by requiring a standardized answer prefix. Our v2 system prompt adds the instruction \texttt{**ANSWER: <number>**}, making extraction a simple regex match on a known prefix. This eliminates false negatives from formatting variation while remaining faithful to the original task.

\paragraph{Code.} The original HumanEval system prompt does not specify how the assistant should structure its code output. In the multi-turn setting, the assistant sometimes omits the function definition entirely (inlining the logic), or assumes input comes from \texttt{stdin} rather than function parameters, because the sharded problem description does not always convey the expected function signature. Our v2 system prompt explicitly requires: (1)~exactly one top-level function definition, (2)~inputs as function arguments rather than \texttt{stdin}, and (3)~a markdown Python code fence for the final answer. We also relaxed the extraction code to handle minor formatting variations (e.g., extra whitespace, trailing comments).

\paragraph{Database.} Similar to code, the original database prompt does not specify an output format, making SQL extraction fragile. Our v2 prompt requires the final query in a \texttt{```sql} markdown code fence, with the constraint that only one such fence should be used for the final answer.

\paragraph{Actions (function calling).} This is the only LiC task that expects multiple output artifacts (a list of function calls). Because shards are revealed incrementally, the assistant sometimes reasonably assumes it does not need to repeat previously established function calls in its final response, but evaluation only considers the last assistant message. Our v2 system prompt adds an explicit instruction to output all function calls needed to fulfill the user's complete request at each turn, ensuring the evaluated response contains the full set.

\subsubsection{False negative identification}
\label{app:false-neg}

Through manual trace analysis, we found that the user simulator model (which converts problem shards into natural-language user messages) sometimes rephrases shards in ways that subtly distort the original meaning, introducing ambiguity, altering constraints, or omitting critical details. While LiC controls for the \emph{sharding} process by verifying that concatenated-shard performance approximates single-turn performance, it does not control for the user simulator's \emph{rephrasing} and \emph{ordering} of those shards.

We account for this with a post-hoc false negative identification procedure applied to all tasks. For each instance where the assistant fails, we collate all user simulator messages and compare them against the original single-turn question. We query GPT-5 (prompt provided in supplemental materials) to determine whether the user simulator's messages, taken as a union, provide sufficient information to solve the original problem.

Instances flagged as insufficient are excluded from accuracy calculations, as the assistant's failure is attributable to the evaluation setup rather than to context pollution or the intervention strategy. Manual analysis of the classifier's outputs finds it to be higher-precision, lower-recall: it conservatively flags only clear cases of information loss or distortion. We consider this acceptable because we prefer to allow some ambiguity to the assistant (counting borderline cases as valid) rather than overly invalidating results.

\subsubsection{Variance exploration on the recommended default}
\label{app:variance}

We do not report error bars on every cell of Table~\ref{tab:main} (see Limitations), but we did spot-check that the headline LiC ordering is not a lucky-run artifact. We re-ran the recommended-default operator (Gated-Reset) two additional times across all four LiC tasks, giving $N{=}3$ per cell when combined with the headline run. We focus the variance budget on a single complete row rather than scattered cells, both so the row is coherent and because Gated-Reset is the operator the paper recommends as the conservative default. Two additional implicit checks come for free from the experiment design: (i)~each of \method's operator variants (Augment, Reset, Gated-Reset, plus their +Memory counterparts) reuses the same upstream analyzer artifact per question, so the four-row ordering across operators is also a check on robustness across different operationalizations of the same analysis step --- and Gated-Reset clears the Full-Context baseline on every LiC task in this row, mirroring the variance result; (ii)~the multi-model results (Table~\ref{tab:multi-model}) replay the same protocol across four model families and recover the same qualitative ordering 15 of 16 cells, providing a cross-model robustness check on the analysis pipeline itself.

\paragraph{Protocol.} We use replay mode (last-turn regeneration with fixed user-sim trajectory loaded from saved S0 baseline traces). This isolates analyzer + assistant variance --- the user simulator's shard-reveal trajectory is held constant across the three runs --- and so likely \emph{underestimates} full end-to-end variance, which would also include user-sim stochasticity. The configuration matches the headline run for each cell: \texttt{context\_edit\_v2} on math/code/database with replay source \texttt{data/baseline\_traces\_v2/\{task\}}, and \texttt{context\_edit\_v2\_accumulate} on actions (which adds the explicit instruction to repeat all function calls in the final turn) with replay source \texttt{outputs/2026-03-16/19-26-46}. All runs use \texttt{gpt-5-mini} as both assistant and analyzer at temperature~1.0.

\paragraph{Results.} Per-run accuracies and aggregate statistics are below (raw $n$, no auxiliary normalization):

\begin{center}
\begin{tabular}{lccccc}
\toprule
\textbf{Task} & \textbf{Headline} & \textbf{Run 2} & \textbf{Run 3} & \textbf{Mean} & \textbf{Std} \\
\midrule
math (n=20)     & 75.0 & 80.0 & 85.0 & \textbf{80.0} & 5.0 \\
code (n=18--19) & 72.2 & 63.2 & 57.9 & \textbf{64.4} & 7.2 \\
database (n=25) & 44.0 & 32.0 & 40.0 & \textbf{38.7} & 6.1 \\
actions (n=25)  & 68.0 & 56.0 & 60.0 & \textbf{61.3} & 6.1 \\
\bottomrule
\end{tabular}
\end{center}

\paragraph{Reading.} Std is in the 5--7pp range. Means cleanly clear the Full-Context baseline (60.0/15.8/4.0/34.8) and the Augment operator without memory (80.0/55.6/32.0/47.8) on every task: $80.0 \pm 5.0$ vs Augment 80.0 (parity, expected --- math is the easy direction); $64.4 \pm 7.2$ vs Augment 55.6; $38.7 \pm 6.1$ vs Augment 32.0; $61.3 \pm 6.1$ vs Augment 47.8. The qualitative ordering --- Gated-Reset substantially above the Full-Context baseline; Reset/Gated-Reset above Augment on context-removal-sensitive tasks (code/database/actions) --- is preserved. Specific point-estimate gap-closure percentages cited in the body (e.g., ``78\% of the gap on code'') should be read as derived from these means.

\paragraph{What this does and does not establish.} This is variance under fixed user-sim trajectories with the analyzer and assistant resampled. It does not estimate variance from re-rolling the user simulator (which would require full simulation, not replay). Reruns of the other ours-rows (Augment, Reset, +Memory) are left to future work; the choice to anchor variance evidence on Gated-Reset reflects that it is the operator recommended for production use in the conclusion.

\paragraph{Note on the \texttt{actions} cell.} An earlier draft reported the actions cell at 73.9\% (= 17/23), reflecting a post-hoc ``common 23-sample'' normalization that excluded samples with user-simulator artifacts, applied non-uniformly across runs. We use raw $n{=}25$ consistently across the three runs reported here for apples-to-apples comparison; this drops the headline value from 73.9\% to 68.0\% but yields a clean per-run comparison.

\subsection{Multi-model evaluation protocol}
\label{app:multi-model-protocol}

Because stronger models solve easier problems regardless of context pollution, we re-ran the subset selection protocol using GPT-5.2 LiC logs (instead of GPT-5-mini), selecting the 20 hardest true-negative instances per task to ensure a challenging set for this model class. Each model serves as both assistant (final-turn regeneration) and analyzer. All models replay conversations originally generated by GPT-5.2; the Baseline reflects how model $X$ handles GPT-5.2's polluted context, while Reset reflects model $X$ given a clean context derived from its own analysis. User and system agents use GPT-4o-mini throughout. We compare Baseline against unconditional Reset (the best-performing strategy in Section~\ref{sec:lic-results}). Denominators vary slightly for code (14--15 valid samples due to evaluation errors); percentages use valid denominators. We additionally evaluate GPT-5, DeepSeek V3.2 (37B active, open-weight), and Qwen 3.5 35B-A3B (3B active, open-weight).

\begin{table}[h]
\centering
\begin{tabular}{llccccc}
\toprule
\textbf{Model} & \textbf{Strategy} & \textbf{Math} & \textbf{Code} & \textbf{Database} & \textbf{Actions} & \textbf{Avg.\ $\Delta$} \\
\midrule
\multirow{2}{*}{GPT-5-mini} & Baseline & 20\% & 10\% & 5\% & 60\% & \\
 & Reset & \textbf{55\%} & \textbf{25\%} & \textbf{40\%} & \textbf{80\%} & +26pp \\
\midrule
\multirow{2}{*}{GPT-5} & Baseline & 40\% & 21\% & 27\% & 85\% & \\
 & Reset & \textbf{70\%} & \textbf{47\%} & \textbf{65\%} & 70\% & +20pp \\
\midrule
\multirow{2}{*}{DeepSeek V3.2} & Baseline & 50\% & 7\% & 10\% & 35\% & \\
 & Reset & \textbf{80\%} & \textbf{62\%} & \textbf{45\%} & \textbf{70\%} & +39pp \\
\midrule
\multirow{2}{*}{Qwen 3.5 35B} & Baseline & 55\% & 21\% & 5\% & 45\% & \\
 & Reset & \textbf{75\%} & \textbf{64\%} & \textbf{75\%} & \textbf{80\%} & +42pp \\
\bottomrule
\end{tabular}
\caption{Multi-model LiC results on a hard subset selected via high baseline failure rate by GPT-5.2. Reset improves accuracy for 15 of 16 model--task pairs across four architectures.}
\label{tab:multi-model}
\end{table}

\subsection{CollabLLM evaluation details}
\label{app:collabllm}

Standard BigCodeBench evaluation uses test-case execution requiring exact function signatures, but the CollabLLM user simulator never conveys these signatures exactly, making pass\_rate systematically zero for all conditions. We develop a conversation-aware LLM judge (GPT-5) that evaluates the assistant's code against the user's actual intent as expressed in the conversation, assessing functional correctness rather than signature conformance. Headline GPT-5-mini numbers (Baseline 40.0\% / 62.5\%, Rewrite 45.0\% / 82.5\% on MATH-Hard / BigCodeBench) are aggregated in the CollabLLM block of Table~\ref{tab:megatable}. Rewrite uses single-pass full-conversation analysis (no structural exclusion).

\subsection{WildChat protocol and turn type definitions}
\label{app:wildchat-turns}

\paragraph{Filtering and protocol.} Following \citet{huang2026context}, we filter WildChat-1M to English, technical conversations of 5--10 rounds. The evaluation has two phases: Phase~1 identifies turns where AO degrades quality; Phase~2 evaluates our strategies on all turns using Huang's pairwise LLM judge prompt, which they validated against human annotators.

\paragraph{Turn type definitions.} WildChat turns are classified into three types based on the relationship between the current user message and prior conversation context: \texttt{new\_ask} (the user introduces a new topic or question unrelated to prior turns), \texttt{feedback} (the user provides feedback on or corrections to the assistant's prior response), and \texttt{no\_feedback} (the user continues the conversation without explicitly referencing the assistant's response, e.g., providing additional information or following up on their own prior message). AO is most harmful on \texttt{no\_feedback} turns, where the assistant's prior work contains relevant state that should not be discarded.

\subsection{Tau2-bench agentic adaptation}
\label{app:agentic-methods}

For tau2-bench, we make three modifications: (1)~the analyzer explicitly tracks environment state (account settings, diagnostic results) revealed through tool calls; (2)~the analysis is reframed as ``strategic reflection,'' identifying unexplored diagnostic avenues rather than just flagging errors; and (3)~resets are gated and capped at 3 per conversation to prevent infinite loops.

\subsection{Tau2-bench results and diagnostic analysis}
\label{app:tau2-diagnostic}

Headline tau2-bench \texttt{telecom\_small} results are reported in Table~\ref{tab:megatable} (tau2 block); the per-trial gpt-5-mini breakdown follows. Per-trial success rates across seeds 42--44 are: Baseline $\{45.0, 55.0, 60.0\}\%$, \method-Gated-Reset $\{40.0, 65.0, 40.0\}\%$, AO $\{0.0\}\%$ (single trial; reward is invariant to seed because no diagnostic state is preserved across an AO context, so re-running with a different seed produces the same 0\%). The table reports the best-of-3 trial for Baseline and Gated-Reset; AO collapses to 0\% on every trial because tool-call results live only in assistant turns. Per-trial variance is wide for both Baseline and Gated-Reset, and the two are statistically within trial noise of each other --- on this stateful regime our method neither helps nor catastrophically hurts the agent's success rate, at $\sim$+20\% cost per task. The headroom for any context-editing intervention is inherently small here: of 11 Baseline failures, only 1 is attributable to context pollution (Table~\ref{tab:tau2-failure-modes}). We hypothesize tau2-bench's customer-service policy structure --- which lays out high-level diagnostic options for the agent to enumerate --- already gives the agent a breadth-first re-entry mechanism around early incorrect paths, partially substituting for the role of an external editor. The decisive contrast in this row is that blanket assistant omission destroys all diagnostic state and so cannot be applied without complete failure, while selective editing remains viable.

Diagnostic analysis of baseline failures on tau2-bench reveals the distribution of failure modes:

\begin{table}[h]
\centering
\begin{tabular}{lrl}
\toprule
\textbf{Category} & \textbf{Count} & \textbf{Description} \\
\midrule
Task too hard & 4 & Agent never tries the right fix \\
Step budget / Hard persona & 4 & Correct approach but runs out of turns \\
Context pollution & 1 & Agent stuck in repetitive loop \\
User sim issue & 1 & User stopped after first message \\
Premature transfer & 1 & Agent transfers instead of troubleshooting \\
\bottomrule
\end{tabular}
\caption{Baseline failure mode breakdown (11 failures). Context pollution is not the dominant failure mode in tau2-bench.}
\label{tab:tau2-failure-modes}
\end{table}

This finding motivated the shift from pollution-detection framing (v9) to strategic-reflection framing (v10), which addresses the dominant failure mode of missing domain knowledge by redirecting the agent toward unexplored diagnostic avenues.


\section{Prompt templates}
\label{app:prompts}

We provide the full text of key prompts used in our system. Template variables are shown in braces (e.g., \texttt{\{system\_message\}}). Task-specific system prompts and the false negative identification prompt are provided in the supplemental materials.

\subsection{Conversation analysis prompts}

\refstepcounter{promptctr}
\paragraph{Prompt \thepromptctr: Subtask 1 --- Task specification extraction (hard attention).}
\label{prompt:query1}
\promptinput{assets/prompts/analyzer_query1_task_spec.txt}

The analyzer sees \emph{only} the system message and user messages; assistant messages are structurally excluded. This is the ``hard attention'' design described in Section~\ref{sec:cognitive-hazard}.

\refstepcounter{promptctr}
\paragraph{Prompt \thepromptctr: Subtask 2 --- Approach evaluation.}
\label{prompt:query2}
\promptinput{assets/prompts/analyzer_query2_compare.txt}

Subtask~2 receives the clean task specification from Subtask~1 and the full conversation (including assistant messages). It evaluates alignment and identifies issues.

\subsection{Context compaction prompt}

\refstepcounter{promptctr}
\paragraph{Prompt \thepromptctr: Context compaction (used by Reset and Gated Reset).}
\label{prompt:compaction}
\promptinput{assets/prompts/context_compaction.txt}

When the edit decision is affirmative, this prompt produces the replacement context that the assistant sees in place of the original conversation history.

\subsection{Memory-based learning prompts}

\refstepcounter{promptctr}
\paragraph{Prompt \thepromptctr: Trajectory reflection (Reflect step).}
\label{prompt:reflect}
\promptinput{assets/prompts/reflect_takeaways.txt}

After each trajectory, this prompt extracts transferable analysis principles. The output is a bullet list of takeaways grounded in the specific trajectory but generalizable across domains.

\refstepcounter{promptctr}
\paragraph{Prompt \thepromptctr: Cheatsheet unification (Unify step).}
\label{prompt:unify}
\promptinput{assets/prompts/unify_takeaways.txt}

The Unify step integrates individual trajectory takeaways into a single coherent cheatsheet, capped at 1500 words to prevent unbounded growth. The cheatsheet is then injected into the analyzer's context via the \texttt{\{memory\_section\}} template variable in Prompts~\ref{prompt:query1} and~\ref{prompt:query2}.


\section{Soft vs.\ hard attention}
\label{app:soft-attention}

Structural exclusion of assistant messages (hard attention) exploits a property of LiC: user messages alone fully specify the task. In realistic settings, the task may depend on prior assistant work. For example, when a user says ``sharpen the phrasing in sentence 2,'' the referent exists only in the assistant's output. \citet{huang2026context} observe similar cases where user messages alone are insufficient for understanding the task.

We introduce a \emph{soft-attention} variant where Subtask~1 sees the full conversation but explicitly reasons through source attribution: listing user-stated requirements with provenance, flagging assistant-originated interpretations as hypotheses, and writing the specification from user-grounded requirements only. We call this ``soft attention'' because the model must rely on transformer attention to selectively weight user-grounded information, rather than having polluting content physically removed from its context.

\begin{table}[h]
\centering
\begin{tabular}{lccc}
\toprule
\textbf{Condition} & \textbf{Math} & \textbf{Code} & \textbf{Database} \\
\midrule
\multicolumn{4}{l}{\emph{Augment}} \\
\quad Soft attention (no memory) & 40\% & 13\% & 31\% \\
\quad Soft attention + memory & 56\% & 14\% & \textbf{46\%} \\
\quad Hard attention (oracle) & \textbf{80\%} & \textbf{40\%} & 54\% \\
\midrule
\multicolumn{4}{l}{\emph{Reset}} \\
\quad Soft attention (no memory) & 42\% & 23\% & 31\% \\
\quad Soft attention + memory & 42\% & 25\% & \textbf{46\%} \\
\quad Hard attention (oracle) & \textbf{67\%} & \textbf{46\%} & 54\% \\
\bottomrule
\end{tabular}
\caption{Soft- vs.\ hard-attention results on held-out test splits. Memory closes up to 65\% of the gap to hard attention on database.}
\label{tab:soft-attention}
\end{table}

Hard attention substantially outperforms soft attention across all tasks (Table~\ref{tab:soft-attention}), confirming that prompting alone cannot substitute for structural exclusion. However, Reset outperforms Augment even with contaminated specs (code: 23\% vs.\ 13\%), confirming that context pollution and spec contamination are separable problems.

Memory-based decontamination (using hard-attention specs as oracle targets to learn decontamination principles via the Dynamic Cheatsheet) shows its strongest effect on database (+15pp), closing 65\% of the gap to hard attention. The learned principles converge on structured heuristics: ``build a user-only spec first, then overlay assistant suggestions as hypotheses.'' Code remains resistant (+1--2pp), suggesting more structurally diverse contamination patterns that a minimal memory approach cannot address.

\section{Hard-attention ablation (Table~\ref{tab:cognitive-hazard})}
\label{app:hard-attention}

\begin{table}[h]
\centering
\begin{tabular}{lccc}
\toprule
\textbf{Configuration} & \textbf{Math} & \textbf{Code} & \textbf{Database} \\
\midrule
Context Edit (Augment) & \textbf{80\%} & \textbf{56\%} & \textbf{32\%} \\
\quad $-$ hard attention & 40\% & 11\% & 8\% \\
\quad\quad $-$ subtask decomposition & 55\% & 21\% & 4\% \\
\midrule
Baseline (no analysis) & 60\% & 16\% & 4\% \\
\bottomrule
\end{tabular}
\caption{Progressive ablation of the analyzer (Augment, no memory). Row~2 lets the spec extractor see assistant messages; Row~3 additionally collapses the two-step analysis into a single pass. Contamination drops accuracy below baseline, and the single-pass variant outperforms the contaminated two-step version because the latter trusts a biased upstream spec as ground truth.}
\label{tab:cognitive-hazard}
\end{table}

The headline numbers in the main paper used structural exclusion: Subtask~1 is given only user messages, and physically cannot see assistant turns. We progressively remove design choices to test whether this is load-bearing or whether prompt-level instructions (``please reason only from user turns'') could substitute. First, we drop hard attention (let Subtask~1 see assistant messages while keeping subtask decomposition); second, we additionally collapse the two-step pipeline into a single pass.

\paragraph{Assistant messages are a cognitive hazard.} Letting the spec extractor see assistant messages collapses performance below baseline on all tasks. Even a model explicitly instructed to act as an independent reviewer anchors on the assistant's reasoning, biasing the extracted task specification.

\paragraph{Contamination compounds across stages.} Counterintuitively, the single-pass variant \emph{outperforms} the contaminated two-step version. When Subtask~1 is contaminated, Subtask~2 treats the biased spec as ground truth and amplifies the upstream error. Decomposition into stages is only beneficial when earlier stages produce clean outputs; otherwise, fewer stages may be safer than more. The takeaway transfers to any multi-stage LLM pipeline in which one model evaluates another's output: structural data-flow control may be needed when prompting alone is insufficient for selective attention.

\section{Memory-based learning details}
\label{app:memory-details}

\begin{table}[h]
\centering
\begin{tabular}{lccc}
\toprule
\textbf{Condition} & \textbf{vs AO} & \textbf{vs FC} & \textbf{$\Delta$ vs AO} \\
\midrule
\multicolumn{4}{l}{\emph{DeepSeek V3.2 (eval split)}} \\
\quad Reset & 72.5\% & 72.5\% & --- \\
\quad\quad + memory & \textbf{86.3\%} & \textbf{80.4\%} & \textbf{+13.8pp} \\
\midrule
\multicolumn{4}{l}{\emph{GPT-5-mini (online)}} \\
\quad Reset & 90.8\% & 86.8\% & --- \\
\quad\quad + memory & 86.7\% & 80.0\% & $-$4.1pp \\
\bottomrule
\end{tabular}
\caption{WildChat memory ablation. Indented \texttt{+ memory} rows add the memory-based learner to the Reset row above. Memory improves DeepSeek's win rate by +13.8pp when there is headroom; GPT-5-mini's near-ceiling baseline leaves no room for improvement. Memory was learned from one-third of WildChat turns and frozen for evaluation. We keep this as a separate panel from the main results because it isolates a cross-model effect that the main \emph{(c)} panel (single-model strategy comparison) does not show.}
\label{tab:wildchat-memory}
\end{table}

We evaluate two learning regimes. \emph{Online}: reflections update the cheatsheet incrementally as problems are solved in batches, so later problems benefit from accumulated principles, a form of continual test-time learning. \emph{Offline}: a set of training trajectories is used to initialize the cheatsheet via a Reflect-then-Unify pipeline (individual reflections followed by a consolidation pass), and the frozen cheatsheet is applied at evaluation time. The offline regime provides a controlled comparison where the only variable is the presence of the cheatsheet. On LiC, we use online learning. On WildChat, we use the offline regime with a train/eval split (one-third training, two-thirds evaluation) to avoid confounding memory quality with ordering effects and to test cross-domain transfer.


\section{Trajectory example: WildChat Maven debugging}
\label{app:trajectories}

We present a representative WildChat conversation where context curation (S2) outperforms both full-context (FC) and assistant-omitted (AO) baselines across all evaluated turns. The conversation (\texttt{300289e3...}, 8 rounds, originally GPT-3.5-turbo) involves iterative debugging of a Maven/Java build in GitHub Codespaces.

\subsection{Conversation summary}

The user runs \texttt{mvn clean install} and encounters a sequence of errors. Early turns establish the environment (GitHub Codespaces) and initial diagnosis (JRE vs.\ JDK). By Turn~8, the user reports \emph{``Still getting: Fatal error compiling: error: invalid target release: 17''} after following the assistant's prior suggestion, a classic \texttt{no\_feedback} turn where the user implicitly references the prior assistant response without providing explicit corrective feedback.

\subsection{Turn 8: three conditions compared}

\paragraph{AO (assistant-omitted).} With all prior assistant messages replaced by \texttt{[Response provided]} placeholders, AO re-suggests \texttt{apt-get install} and checking \texttt{java -version}, repeating steps the user has already performed. AO cannot connect ``still getting'' to what was already tried.

\paragraph{FC (full context).} FC sees the full history and includes \texttt{update-alternatives} commands, but still partially repeats earlier suggestions. The model anchors on the same diagnostic approach (checking \texttt{java -version}) rather than recognizing that the user already tried the basic install without resolution.

\paragraph{S2 (context curation).} The analyzer identifies that the prior assistant's diagnostic reasoning was flawed (relying on \texttt{java -version} text parsing rather than checking \texttt{javac} directly) and that the conversation was stuck repeating an insufficient fix. The rewritten context preserves the correct diagnosis (Maven needs a JDK) while removing the flawed reasoning. S2's response:
\begin{itemize}
    \item Correctly diagnoses why the error persists: Maven may use a different Java than what \texttt{java -version} reports
    \item Provides non-redundant steps (does not re-suggest things already tried)
    \item Includes solutions FC and AO both missed: \texttt{devcontainer.json} configuration, the \texttt{\textless release\textgreater 17\textless /release\textgreater} tag in \texttt{pom.xml}
\end{itemize}

\paragraph{Analyzer output (abbreviated).}
\begin{quote}
\emph{Issue:} The assistant said that seeing ``OpenJDK Runtime Environment'' in \texttt{java -version} means the system has a JRE rather than a JDK. That is not reliable: \texttt{java -version} output often contains ``Runtime Environment'' even when a JDK is installed. The correct checks are presence of \texttt{javac} and the value of \texttt{JAVA\_HOME}.

\emph{Aligned:} Maven's ``No compiler is provided'' message means Maven cannot find \texttt{javac} and needs a JDK. Installing a JDK inside Codespaces is a valid remediation path.
\end{quote}

\subsection{Judge evaluations}

All three pairwise comparisons were judged by GPT-5-mini:

\begin{itemize}
    \item \textbf{FC vs.\ AO:} FC wins. ``Both responses correctly diagnose the problem. [FC] is slightly clearer and more prescriptive.''
    \item \textbf{AO vs.\ S2:} S2 wins. ``[S2] is substantially more complete and actionable.''
    \item \textbf{FC vs.\ S2:} S2 wins. ``[S2] is significantly more complete and actionable. It gives precise diagnostic commands, exact installation commands, environment variable exports, pom.xml changes, devcontainer recommendations, and verification steps.''
\end{itemize}

\subsection{Conversation-level results}

Across all 7 evaluated turns, S2 wins every pairwise comparison against both baselines (Table~\ref{tab:traj-turns}). On turns where AO beats FC (turns 2, 10, 12, 14), context pollution hurts more than context loss, yet S2 still wins by removing the pollution while preserving useful context.

\begin{table}[h]
\centering
\begin{tabular}{clccc}
\toprule
\textbf{Turn} & \textbf{User message} & \textbf{FC vs AO} & \textbf{S2 vs AO} & \textbf{S2 vs FC} \\
\midrule
2 & ``Why do I keep getting error\ldots'' & AO & S2 & S2 \\
4 & ``I am using Codespaces'' & FC & S2 & S2 \\
6 & ``Now I get: [new error]'' & FC & S2 & S2 \\
8 & ``Still getting [same error]'' & FC & S2 & S2 \\
10 & ``I do not have a bashrc file'' & AO & S2 & S2 \\
12 & [Shares GitHub Actions YAML] & AO & S2 & S2 \\
14 & ``No WAR file generated'' + pom.xml & AO & S2 & S2 \\
\bottomrule
\end{tabular}
\caption{Per-turn winners for the WildChat Maven debugging conversation. S2 wins all 14 pairwise comparisons. On context-dependent turns (4, 6, 8), AO fails because it cannot connect the user's references to prior discussion.}
\label{tab:traj-turns}
\end{table}


\section{Additional related work}
\label{app:additional-related}

\paragraph{Multi-turn degradation and context pollution (extended).}
\citet{laban2025lost} introduced the Lost in Conversation (LiC) framework, showing that multi-turn performance drops 39\% on average across 15 LLMs, driven by increased unreliability (+112\%) in addition to reduced capability ($-$16\%). \citet{huang2026context} termed this \emph{context pollution} and showed that simply omitting all assistant messages recovers much of the lost performance, establishing that the assistant's own outputs are the primary source of degradation. \citet{khalid2025ergo} proposed ERGO, which monitors token-level entropy and triggers context resets via an LLM rewrite of prior user messages. Both approaches treat context as uniformly harmful: this is effective only for self-contained tasks where the model can re-derive answers from user messages alone, a setting LiC was explicitly designed to test (each instance is a single-turn question sharded across turns). This result has been read as evidence that multi-turn degradation is largely cosmetic, but on referential tasks (WildChat) and especially stateful agentic tasks (tau2-bench) blanket omission collapses from near-oracle to zero, which shows that the underlying pollution still exists; omission simply happens not to help when assistant turns carry referential or stateful content. We build on their diagnostic finding while treating context as heterogeneous, and show that fine-grained editing closes most of the remaining LiC gap without collapsing the conversation to a single turn while remaining viable in the agentic setting. Note that ERGO is not a design oracle: its LLM rewrite step can paraphrase or compress the user's exact phrasing, which is why it underperforms simple concatenation on database and actions (Table~\ref{tab:main}).

\paragraph{Test-time scaling.}
Methods such as chain-of-thought~\citep{wei2022chain}, self-consistency~\citep{wang2023selfconsistency}, and structured search~\citep{yao2024tree} scale inference compute within a single generation. Our work extends test-time scaling to the multi-turn setting: rather than searching over possible outputs, we edit the \emph{input context} presented to the model across turns, investing additional LLM calls to improve the quality of information the model reasons over.

\paragraph{Memory-based self-improvement.}
A growing line of work uses persistent memory to improve LLM capabilities without parameter updates. General-purpose memory systems store and retrieve knowledge for personalization and continuity~\citep{xu2025amem,zhou2025mem1}. Agentic memory systems go further, accumulating procedural knowledge and task-specific heuristics that improve agent performance over time: the Dynamic Cheatsheet~\citep{suzgun2025cheatsheet} maintains a mutable text document updated across problem instances; Voyager~\citep{wang2023voyager} builds a skill library for open-ended exploration; and several recent systems learn reusable reasoning strategies~\citep{ouyang2026reasoningbank,ho2025arcmemo,fang2026memp,wu2025evolver,zhang2026ace}. A separate thread optimizes the memory system itself rather than the downstream task~\citep{zhang2025memevolve,xiong2026alma}. We adapt the Dynamic Cheatsheet for context curation, targeting accumulated analysis principles at the analyzer's executive function rather than the task-solving assistant.

\section{Extended agentic context management discussion}
\label{app:agentic-context}

THREAD~\citep{schroeder2025thread} spawns sub-agents that report summarized results back to a main thread, reducing context accumulation. Recursive Language Models~\citep{zhang2025rlm} treat context as a variable manipulated via code operations in a REPL, enabling a controller LLM to programmatically decompose and recursively query sub-calls. Context Folding~\citep{sun2025contextfolding} uses reinforcement learning (FoldGRPO) to train agents that branch into sub-trajectories and collapse intermediate steps into concise summaries, maintaining active context 10$\times$ smaller than ReAct baselines.

These methods share a common goal, \emph{compaction}: fitting more reasoning into a bounded context window. Our work addresses a fundamentally different problem: \emph{context pollution}, the assistant's own prior outputs actively mislead future generation, even well within the context window. Compaction methods preserve information while reducing tokens; we selectively \emph{remove} information that is harmful. The two goals are complementary but require different analytical capabilities.

U-Fold~\citep{su2026ufold} targets multi-turn conversation with a two-stage pipeline: intent summarization (creating a TODO-list from conversation history) and data extraction (retaining only relevant tool responses). While structurally similar to our two-query decomposition, U-Fold addresses redundancy in tool responses, not self-induced pollution from the assistant's own reasoning.

MemoBrain~\citep{qian2026memobrain} formalizes ``executive memory'' as an online mechanism that constructs a dependency-aware graph of reasoning trajectories and exerts cognitive control over working context. Like our work, it draws on the executive function framework and controls information flow to the reasoning agent. However, MemoBrain's executive control is triggered by approaching budget limits and its operations (folding sub-trajectories into thought units, flushing low-salience nodes) serve long-horizon compaction. Our executive function targets a different deficit: the inability to \emph{disregard} prior reasoning that has been invalidated, even when the model has ample context budget.

\section{Executive function as a design pattern}
\label{app:exec-function-pattern}

Our approach instantiates a general pattern: \emph{externalized executive function} via an agent that schedules LLM calls with custom context windows. Different subtasks of self-reflection require different context: intent extraction must be shielded from contamination, approach evaluation needs both clean intent and the full history, strategic reflection identifies unexplored alternatives. This pattern generalizes: any setting where an LLM reasons about potentially misleading context (code review, fact-checking) could benefit from structured context control. The executive function framing provides the design vocabulary: what should each reasoning step attend to, and what should it be shielded from?

We do not claim the LLM is self-aware in any phenomenological sense. Rather, the system designer's knowledge of context pollution as a systematic failure mode is encoded into the agent architecture, analogous to using an external checklist to compensate for known cognitive biases.

\section{Limitations}
\label{app:limitations}

Several limitations qualify the results in this paper. \emph{Simulated users.} LiC, CollabLLM, and tau2-bench use simulated users; behavior with real users may differ in ways the simulator does not capture. \emph{Replay protocol for multi-model results.} The multi-model evaluation (Section~\ref{sec:multi-model}) replays conversations originally generated by GPT-5.2 across other models, isolating the effect of context curation given a fixed history but not testing each model's own multi-turn generation dynamics; the resulting gains should be read as ``how much does this method recover from a polluted history,'' not ``how much does it improve this model's native multi-turn behavior.'' \emph{Sample size and statistical reporting.} Per-cell sample sizes are modest and we do not report paired significance tests on the main tables; the headline gaps should therefore be read as suggestive of the qualitative trends Figure~\ref{fig:story} summarizes rather than as point-calibrated effect sizes. \emph{Tau2-bench scope.} Tau2-bench results use only the \texttt{telecom\_small} subset and are framed as a preliminary exploration of the agentic regime; the load-bearing finding is the \emph{collapse} of AO when tool-call results live only in assistant turns, not the absolute success rate of \method. \emph{Untested settings.} Long-horizon tasks (multi-file coding, iterative design) and broader agentic domains remain untested.


\newpage
\input{checklist.tex}


\end{document}
